% ==================================================================
% id:       app-007
% topic:    שימושי הנגזרת
% subtopic: בעיות קיצון
% solution: full
% ==================================================================
מבין כל המלבנים שהיקפם \(P\), מצאו את מידות המלבן ששטחו המקסימלי.
\begin{solbox}
נסמן את אורך של המלבן ב-
$x$
ורוחב של המלבן ב-
$y$.
מכיוון שההיקף הוא
$P$
מתקיים
$$,2x+2y=P$$
מכאן
$y=\frac{P}{2}-x$.
מכאן השטח של מלבנים כאלה נתון ע"י
$$.xy=x\left(\frac{P}{2}-x\right)=f(x)$$
עלינו למקסם את הפונקציה
$f$
בתחום
$0\leq x\leq\frac{P}{2}$
(שימו לב כי
$y=\frac{P}{2}-x\geq0$,
לכן
$x\leq\frac{P}{2}$).
נגזור ונשווה את הנגזרת ל-
$0$:
$$,0=f'(x)=\frac{P}{2}-2x$$
קיבלנו כי
$x=\frac{P}{4}$.
מתקיים
$f(\tfrac{P}{4})=\frac{P^2}{16}$.
נבדוק את ערכיה של
$f$
בקצוות הקטע:
$$.f(0)=0,\ f(\tfrac{P}{2})=0$$
אנו רואים כי
$$,f(\tfrac{P}{4})>f(0)=0,\ f(\tfrac{P}{4})>f(\tfrac{P}{2})=0$$
נסיק כי ערכה המירבי של
$f$
בקטע
$[0,\tfrac{P}{2}]$
מתקבל בנקודה
$x=\frac{P}{4}$.
בנקודה זו
$y=\frac{P}{2}-\frac{P}{4}=\frac{P}{4}$.
שימו לב כי שטח מירבי מתקבל עבור ריבוע.
\end{solbox}
